\section{Introducción}
\label{sec:introduccion}

El artículo \emph{Cross-lingual distillation for domain knowledge transfer with sentence transformers}, de Piperno, Bacco, Dell'Orletta, Merone y Pecchia, publicado en 2025 en el volumen 311 de \emph{Knowledge-Based Systems}, aborda la subrepresentación de las lenguas distintas del inglés en el procesamiento de lenguaje natural y elige el dominio biomédico \cite{piperno2025}. Su propuesta traduce textos médicos ingleses al español, al francés y al alemán con mBART, alinea las representaciones de un profesor entrenado en inglés con las de un estudiante mBERT mediante \emph{Multiple Negatives Ranking Loss} y ajusta el estudiante con etiquetas en la lengua destino usando entre el 10 y el 100\,\% de los datos. El máximo reportado es de 12,89 puntos de F1 con el 10\,\%; el promedio cae de 7,95 a 0,45 puntos al pasar al 100\,\%.

El trabajo se inscribe en el seminario «Más allá del modelo», cuyas pautas piden contar el artículo como una historia y separar lo que reportan los autores de la valoración del equipo \cite{moreno2026}. La lectura sigue esa indicación: no resume sección por sección, reconstruye el argumento y lo examina en los puntos donde la evidencia no alcanza.

Nuestra tesis es que la destilación es una técnica de eficiencia de datos, no una garantía de equidad ni de validez clínica. El artículo abre con un marco de inclusión tecnológica y cierra con la promesa de hacer accesible la información biomédica en varias lenguas, pero su diseño mide una sola cosa: la exactitud de un clasificador temático en tres lenguas europeas. Tres observaciones sostienen la tesis. La ganancia se concentra en el régimen de pocos datos y se desvanece cuando hay anotación suficiente. El término \emph{unbiased dataset} se refiere a la ausencia de solapamiento léxico entre el entrenamiento traducido y la prueba, no a sesgo social. La evaluación se hace sobre manuales médicos, no sobre texto clínico real, así que la validez externa hacia registros hospitalarios o hacia el español de Colombia queda sin comprobar.

